\section{Discussion}

\textbf{The heterophilic gains are large and outside noise.}~~The +4.1\% on Chameleon and +3.7\% on Squirrel over SGFormer each exceed one standard deviation of both methods. PCGT also beats every heterophily-specific baseline by wide margins: GloGNN by +8.8\%/+9.8\% and H$_2$GCN by +10.9\%/+10.4\% on Chameleon/Squirrel respectively. The partition-conditioned attention mechanism appears to capture community-level structure that flat global attention (SGFormer) and fixed coefficient matrices (GloGNN) miss.

\textbf{Homophilic improvements: hybrid benefit, not attention alone.}~~Paired $t$-tests (10 runs) give $p$=0.048 on CiteSeer, $p$=0.027 on PubMed, and $p$=0.004 on Coauthor-Physics. Cora is a wash ($p$=0.88). With $\lambda_{\text{gw}}=0.8$ these datasets receive 80\% of their signal from the GCN branch, so the improvements likely reflect regularization from the partition attention rather than the attention mechanism outperforming local aggregation. PCGT's variance is lower across the board, consistent with this regularization interpretation.

\textbf{$\beta$ reflects feature--context complementarity, not homophily.}~~Multi-run verification across all 11 datasets shows that the learned self-connection weight $\beta$ does not track homophily ($r = 0.006$, $p = 0.99$). Instead, $\beta$ appears to reflect the informativeness of a node's own features relative to the attention-weighted context (noting that $n=11$ limits the statistical power of the correlation test). Nine of 11 datasets have $\beta > 0$, including the heterophilic Chameleon ($\beta = +3.08 \pm 0.02$, 5 runs) and Squirrel ($\beta = +3.53 \pm 0.04$, 5 runs), which have the highest-dimensional features (2325 and 2089 dims respectively). On Film ($\beta = -0.99 \pm 1.43$, negative in 8/10 splits) and Deezer ($\beta = -0.57 \pm 0.25$, negative in 5/5 runs), the model suppresses self-features. A standard fixed residual connection ($\beta \equiv 1$) cannot capture this dataset-specific behavior.

\textbf{When partitioning helps (and when it does not).}~~Computing Newman modularity $Q$ of the $K\!=\!10$ METIS partitions reveals that Chameleon ($Q=0.58$, edge-cut ratio 0.22) has tighter clusters than Squirrel ($Q=0.38$, edge-cut ratio 0.32). Paradoxically, the ablation gap from switching to random partitions is \emph{larger} on Squirrel ($-$0.68 vs.\ $-$0.26 on Chameleon): because Squirrel's graph is denser (avg.\ degree 42 vs.\ 20) and less modular, random partitions cut far more inter-community edges and destroy informative structure. On Deezer ($h \approx 0.5$), where edges are nearly label-random, reducing $\lambda_{\text{gw}}$ to 0.5 gives the partition attention more weight and yields 67.2---only +0.1\% over SGFormer, a gap that is within the error bars and not statistically significant.

\textbf{Why PCGT beats GloGNN by so much.}~~GloGNN~\citep{li2022glognn} learns a dense $N \times N$ coefficient matrix and achieves 40.2\% on Chameleon and 35.7\% on Squirrel (from SGFormer Table~2). PCGT's +8.8\% and +9.8\% gains suggest that partition-conditioned attention with learned representatives captures inter-community relationships more effectively than a closed-form global matrix---though we note GloGNN was not designed for these specific filtered splits.
